\documentclass[11pt]{article}
\usepackage[T1]{fontenc}
\usepackage{lmodern}
\usepackage[margin=1in]{geometry}
\usepackage{amsmath,amssymb,amsthm,mathtools}
\usepackage{booktabs,tabularx,longtable,array,microtype}
\usepackage[hidelinks]{hyperref}
\newtheorem{theorem}{Theorem}[section]
\newtheorem{lemma}[theorem]{Lemma}
\newtheorem{proposition}[theorem]{Proposition}
\theoremstyle{definition}
\newtheorem{example}[theorem]{Example}
\newtheorem{definition}[theorem]{Definition}
\theoremstyle{remark}
\newtheorem{remark}[theorem]{Remark}
\newcommand{\E}{\mathbb E}
\newcommand{\Ent}{\mathsf H}
\newcommand{\KL}{\operatorname{KL}}
\newcommand{\kl}{\operatorname{kl}}
\newcommand{\OR}{\operatorname{OR}}
\newcommand{\Comp}{\operatorname{comp}}
\newcommand{\cstar}{c_\star}
\newcommand{\Glg}{\Gamma_{\mathrm{lg}}}
\newcommand{\Gfr}{\Gamma_{\mathrm{fr}}}
\newcommand{\er}{e_{\mathrm{quarter}}}
\newcolumntype{Y}{>{\raggedright\arraybackslash}X}
\numberwithin{equation}{section}
\setlength{\emergencystretch}{2em}
\setlength{\parskip}{0.3em}
\title{When Pooling Cannot Save a Test\\[0.4em]
\large A guided proof for adaptive group testing with a known defective count}
\author{}
\date{October 25, 2026}
\begin{document}
\maketitle

\begin{abstract}
Suppose we know that exactly \(d\) of \(n\) items are defective. A test
of a group of items tells us only whether the group contains at least
one defective. We may choose each new group after seeing the earlier
answers. Can this freedom guarantee fewer tests than examining items
one at a time?

This exposition proves that the answer is no when
\(n\ge2^{31}\) and \(d/n\ge369056/1000000\): the optimal worst-case
number is exactly \(n-1\). We begin with small examples, explain
decision trees, and develop all the probability tools used in the proof.
In particular, entropy, relative entropy, conditioning on an exact
count, and an artificial comparison experiment are introduced from
finite-sum definitions. A structural argument about legal pooled tests
then connects a probability estimate near density one quarter to the
stated theorem. Every mathematical ingredient needed for the conclusion
is proved here. The large numerical constants are accompanied by exact
certificates, with the longer arithmetic placed in an appendix.
\end{abstract}

\tableofcontents
\newpage

\section{The problem: what can a pooled test tell us?}\label{sec:problem}

\subsection{Items, defective sets, and answers}

Imagine five labeled items, numbered \(1,2,3,4,5\). Each item is either
defective or nondefective. The set of defective labels might be
\[
 D=\{2,5\}.
\]
The label \emph{defective} simply names one of two possible
statuses. No statistical model is initially imposed on those statuses.
There is a fixed unknown set, and we must identify it exactly.

A \emph{pool} is a subset of the labels. Testing the pool
\(S=\{1,2,3\}\) gives a positive answer for this \(D\), because label
\(2\) lies in both sets. Testing \(\{1,3\}\) gives a negative answer,
because neither label is defective. A negative answer certifies that
every item in the tested pool is nondefective. A positive answer says
only that at least one item in it is defective. It does not say which
one, or how many.

Write \([n]=\{1,\ldots,n\}\). Formally, the answer to pool
\(S\subseteq[n]\) on defective set \(D\subseteq[n]\) is
\[
 \OR_D(S)=
 \begin{cases}
 0,&S\cap D=\varnothing,\\
 1,&S\cap D\ne\varnothing.
 \end{cases}
\]
This is called a pooled-\(\mathrm{OR}\) test: the answer is the logical
OR of the individual defective indicators in the pool. We call \(0\)
negative and \(1\) positive.

Throughout this document the tests are noiseless, so an answer is
always correct in an actual execution of the testing algorithm.
The algorithm is deterministic: a given history of answers determines
its next query. It is fully adaptive: that query may depend on the
entire history, and may be any subset of the labels.

\subsection{An adaptive strategy must handle all its branches}

\begin{example}\label{ex:four}
There is exactly one defective among four items. First test
\(\{1,2\}\).
If the answer is positive, test \(\{1\}\): the defective is \(1\)
on a positive answer and \(2\) on a negative one.
If the first answer is negative, test \(\{3\}\): the defective is
\(3\) on a positive answer and \(4\) on a negative one.
Thus two tests suffice on every possible input.
\end{example}

The second test depends on the first answer. If the actual defective
is \(3\), only the negative-then-positive branch is executed.
Nevertheless, a correct algorithm must specify successful behavior
on the other three inputs as well. Finding a short explanation for
one particular defective set is easier than designing a strategy
that identifies every possible defective set within a uniform test
limit. Our lower bound concerns the latter requirement.

\subsection{What is known in advance?}

The labels and the exact integer \(d\) are known. The unknown input is
one of the \(d\)-element subsets of \([n]\). We write
\(\binom{[n]}d\) for this family of sets, and \(\binom nd\) for its
number of members. The promise is always \(0<d<n\).
An algorithm may use the exact count at any stage. Once it has
identified all \(d\) defectives, for instance, it knows every other
label is nondefective without testing them.

Knowing the count is different from knowing only an approximate
defective fraction. If four tested items have revealed one defective
and the total is known to be two among five, then the fifth item is
defective. Without the exact total, that conclusion would be unjustified.

\begin{definition}
The \emph{worst-case number of tests} of an algorithm is the largest
number it uses on any \(D\in\binom{[n]}d\).
The number \(M(d,n)\) is the minimum of this worst-case number over
all correct deterministic adaptive pooled-\(\mathrm{OR}\) algorithms.
\end{definition}

The order matters: first take the worst input for a fixed algorithm,
then choose the best algorithm. Different algorithms may have
different worst inputs.

\begin{proposition}\label{prop:individual}
For \(0<d<n\), \(M(d,n)\le n-1\).
\end{proposition}
\begin{proof}
Test labels \(1,\ldots,n-1\) individually. If they contain \(r\)
defectives, the last item must have defective indicator \(d-r\).
For a valid input this indicator is either zero or one, so its
status is determined. This is a correct algorithm using at most
\(n-1\) tests.
\end{proof}

Sometimes this strategy stops earlier, but its bound on all inputs
is \(n-1\). A pooled strategy would improve on it in the worst case
only if it guaranteed at most \(n-2\) tests for \emph{every} valid input.

\subsection{The theorem we will prove}

\begin{theorem}[Individual testing in a dense regime]\label{thm:main}
Let \(\cstar=369056/1000000\). For every pair of integers \(n,d\)
satisfying
\[
 n\ge2^{31},\qquad 0<d<n,\qquad \frac dn\ge\cstar,
\]
one has
\[
 \boxed{M(d,n)=n-1.}
\]
\end{theorem}

In ordinary language: for every population in the stated size range,
if at least the fraction \(369056/1000000\) is defective, no legal
adaptive pooling strategy can guarantee a saving over individual
testing. Some inputs might be identified very quickly; the theorem
says that every correct strategy has at least one input requiring
\(n-1\) or more tests.

Proposition~\ref{prop:individual} already gives the upper bound.
The work lies in the opposite direction, \(M(d,n)\ge n-1\).
Thus this is a lower-bound proof for the number of tests. It must
account for every possible strategy, including strategies whose pools
have complicated sizes and overlaps.

The number \(2^{31}=2147483648\) is an explicit sufficient cutoff in
the proof. We do not claim it is necessary. The theorem has no further
exceptions within its range, but makes no corresponding general
assertion for smaller populations.

\section{Decision trees and the first counting bound}\label{sec:trees}

\subsection{Putting a strategy on paper}

A deterministic adaptive algorithm can be drawn as a rooted binary
tree. An internal node \(v\) contains a pool \(S_v\). Its two outgoing
edges have labels \(0\) and \(1\). A test answer tells us which edge
to take. A leaf \(\ell\) is a stopping point with an output defective
set. Correctness means that every valid input reaches a leaf whose
output is that input.

A \emph{history} or \emph{prefix transcript} \(\tau\) is the finite
string of answers seen so far. In a fixed tree it determines a node,
as long as execution has not already stopped. A \emph{terminal
transcript} \(y\) is the full answer string leading to a leaf.
We distinguish the string \(y\) from the leaf \(\ell\), although the
tree gives a one-to-one correspondence between them.
The depth of a leaf is its number of preceding tests.

Here is Example~\ref{ex:four} as a table:
\begin{center}
\begin{tabular}{cll}
\toprule
Defective set&Queries along its path&Answers\\
\midrule
\(\{1\}\)&\(\{1,2\}\), then \(\{1\}\)&\(11\)\\
\(\{2\}\)&\(\{1,2\}\), then \(\{1\}\)&\(10\)\\
\(\{3\}\)&\(\{1,2\}\), then \(\{3\}\)&\(01\)\\
\(\{4\}\)&\(\{1,2\}\), then \(\{3\}\)&\(00\)\\
\bottomrule
\end{tabular}
\end{center}
The same first query appears in all rows. The second query is a
function of the first answer, not a function of an unknown input
that the algorithm has somehow been told.

We only need finite trees. A strategy with a finite worst-case test
bound can be restricted to its valid-input paths, giving a finite
tree. An algorithm without a finite worst-case bound cannot improve
the finite strategy of Proposition~\ref{prop:individual}.

\subsection{Counting leaves}

\begin{proposition}\label{prop:counting}
Every correct tree for \(\binom{[n]}d\) has worst-case depth at least
\(\left\lceil\log_2\binom nd\right\rceil\).
\end{proposition}
\begin{proof}
Two different valid inputs cannot terminate at the same leaf, because
the leaf has only one output. Thus at least \(\binom nd\) valid
leaves are needed.
A binary tree of depth at most \(s\) has at most \(2^s\) leaves.
This follows by induction: at depth zero there is at most one;
each of two child subtrees has at most \(2^{s-1}\) leaves.
Therefore \(2^s\ge\binom nd\). Taking logarithms and using that
\(s\) is an integer gives the assertion.
\end{proof}

This bound uses only the two possible answers per test. It would hold
even if arbitrary yes/no questions about the input were allowed.
Such questions can be more flexible than pooled tests.

For instance, when \(d=n-1\), there are only \(n\) possibilities:
the locations of the one nondefective item. Counting gives only
\(\lceil\log_2 n\rceil\). Yet every pool of size at least two is
always positive, so it reveals nothing. We will prove in
Section~\ref{sec:structural} that the actual optimum is \(n-1\).
This shows why counting alone cannot establish the desired theorem.

\subsection{A pooled test cannot make an arbitrary split}

Suppose \(n=4,d=2\), giving six candidate inputs.
A negative answer to \(S\) leaves exactly the candidates contained
in \([4]\setminus S\). If \(|S|=k\), their number is
\(\binom{4-k}{2}\), interpreted as zero if \(4-k<2\).
Depending on \(k\), this number is \(6,3,1,\) or \(0\), never \(2\).

An unrestricted binary question could ask whether the input belongs
to a chosen two-element collection of candidates. A pooled test
cannot produce that negative branch on this complete family.
Later the candidate family is more complicated, but an answer must
still be determined by whether the pool intersects the defective set.
Replacing a legal query by a conveniently balanced arbitrary binary
partition would change the problem.

\section{How the proof fits together}\label{sec:roadmap}

The proof has two main parts. One is probabilistic and measures what
a legal tree can accomplish near defective density \(1/4\).
The other is structural and relates optimal test counts on different
population sizes. Neither part alone proves the theorem.

For the probabilistic part we choose a uniformly random valid
defective set solely as a tool for analyzing a fixed deterministic
tree. Let
\[
 H=\log_2\binom nd,\qquad L=\text{its expected number of tests}.
\]
We will explain entropy and prove \(L\ge H\). The difference
\(\Gamma=L-H\) measures how much expected test count exceeds this
basic information requirement.

The important estimate is stronger:
\begin{equation}\label{eq:roadmapWT}
 6615(L-H)\ge n-H-67229K,\qquad K=nh(d/n)-H,
\end{equation}
for \(249/1000\le d/n\le251/1000\).
Here \(h\) is the entropy of one binary outcome, defined below.
The quantity \(K\) is a cost for comparing exact-count inputs with
independent inputs. We will prove \(0\le K\le\log_2(n+1)\).
Thus it is a logarithmic correction, rather than a new linear term.

Why introduce independent inputs? Independence makes probabilities
of all-good pools easy to calculate. But the real inputs have an exact
count and are not independent. Relative entropy will keep track of
the difference over an entire tree.
We also introduce a second comparison experiment that answers large
pools by fair coin flips. This artificial experiment is not a testing
algorithm for the original problem. It leaves certain input coordinates
independent after a transcript, allowing us to calculate how much
uncertainty remains.

The structural part proves a recurrence for \(M\). In a hypothetical
instance where some strategy saves a test, the recurrence forces
smaller instances with increasingly strong savings. Following this
implication eventually reaches a density just below \(1/4\), where
\eqref{eq:roadmapWT} applies. The forced saving is then too large to
be compatible with that inequality.

This explains why an estimate at density \(1/4\) matters for a theorem
whose initial density is at least \(\cstar\). The recurrence connects
the two regimes; the initial problem is not assumed to have density
\(1/4\).

\begin{center}
\begin{tabular}{c}
\fbox{\parbox{0.82\linewidth}{\centering
Finite probability tools + restrictions of pooled-OR queries}}\\[0.4em]
\(\Downarrow\)\\[0.4em]
\fbox{\parbox{0.82\linewidth}{\centering
A whole-tree inequality near density \(1/4\)}}\\[0.4em]
\(\Downarrow\quad\) together with a legal-tree recurrence\\[0.4em]
\fbox{\parbox{0.82\linewidth}{\centering
A supposed saving forces an impossible smaller instance}}\\[0.4em]
\(\Downarrow\)\\[0.4em]
\fbox{\parbox{0.82\linewidth}{\centering
The exact lower bound \(M(d,n)\ge n-1\)}}
\end{tabular}
\end{center}

All probability spaces below are finite. Randomness is introduced to
prove a lower bound, not to allow randomized testing strategies.
Every expectation relevant to the real algorithm averages over valid
exact-\(d\) inputs.

The notation is introduced gradually. A fixed input is \(D\), a
random input is \(X\), a query is \(S_v\), a node is \(v\), a prefix
transcript is \(\tau\), a terminal transcript is \(y\), and a leaf
is \(\ell\). The random terminal transcript is \(Y\).
We write \(P_X\) for an input law, \(P_Y\) for a transcript law,
and \(P_{X,Y}\) for their joint law. This subscript convention also
applies to the comparison laws.
A full notation guide appears in Appendix~\ref{app:notation};
each symbol is explained in the text before it is used.

\section{Entropy: uncertainty measured by a finite sum}\label{sec:entropytools}

\subsection{Why logarithms appear}

An outcome of probability \(1/8\) has the same probability as three
specified successive fair binary choices. This motivates assigning
it information content \(-\log_2(1/8)=3\). More generally, assign
an outcome of probability \(u>0\) information content \(-\log_2u\).
The identity
\[
 -\log_2(uv)=-\log_2u-\log_2v
\]
lets us add information along conditional choices. This logarithmic
identity, rather than a coding theorem, is what we will use.
From now on \(\log\) means \(\log_2\), and \(\ln\) means natural
logarithm. The unit is a bit.

\begin{definition}
For a law \(\mu\) on a finite set \(\Omega\), its \emph{entropy} is
\[
 \Ent(\mu)=-\sum_{z\in\Omega}\mu(z)\log\mu(z).
\]
We use \(0\log0=0\), as justified by the limit when the probability
tends to zero.
\end{definition}

Entropy is the average information content under \(\mu\). It is
nonnegative because \(-\log\mu(z)\ge0\) for every outcome of positive
probability. For a uniform law on \(N\) outcomes,
\[
 \Ent(\mu)=-N\frac1N\log\frac1N=\log N.
\]
A certain outcome has entropy zero. Information content of an
individual outcome and entropy of a law are distinct: the latter
averages the former.

\subsection{One binary outcome}

The entropy of two answers, with probabilities \(u\) and \(1-u\), is
\[
 h(u)=-u\log u-(1-u)\log(1-u),\qquad 0\le u\le1.
\]
It equals one at \(u=1/2\), and zero at \(u=0\) or \(1\).
For example, direct substitution gives
\(h(1/4)=2-\tfrac34\log3\).
Calculus quantifies the fact that biased answers have less entropy:
\[
 h'(u)=\log\frac{1-u}{u},\qquad
 h''(u)=-\frac1{(\ln2)u(1-u)}<0\quad(0<u<1).
\]
The derivative is positive below \(1/2\) and negative above it.
Together with the endpoint values, this proves
\begin{equation}\label{eq:hbound}
 0\le h(u)\le1,\qquad h(u)=1\ \Longleftrightarrow\ u=1/2.
\end{equation}

\subsection{Conditional entropy and its chain rule}

For finite random variables \(U,V\), their \emph{joint law} records
the probabilities of pairs:
\(\mu_{U,V}(u,v)=\Pr(U=u,V=v)\).
Their \emph{marginal laws} record the probabilities of just one
variable; for example,
\(\mu_V(v)=\sum_u\mu_{U,V}(u,v)\).
When \(\mu_V(v)>0\), the conditional law is
\[
 \mu_{U\mid V=v}(u)=\frac{\mu_{U,V}(u,v)}{\mu_V(v)}.
\]
Conditioning values of zero probability are omitted from conditional
averages. With this notation, define
\[
 \Ent_\mu(U\mid V)
 =\sum_{v:\mu_V(v)>0}\mu_V(v)\Ent(\mu_{U\mid V=v}).
\]
This is an average of entropies of conditional laws, weighted by how
often each conditioning value occurs.

\begin{lemma}[Entropy chain rule]\label{lem:entropychain}
\(\Ent(\mu_{U,V})=\Ent(\mu_V)+\Ent_\mu(U\mid V)\).
\end{lemma}
\begin{proof}
For every pair of positive joint probability,
\(\mu_{U,V}(u,v)=\mu_V(v)\mu_{U\mid V=v}(u)\).
Take negative logarithms and average:
\begin{align*}
 -\sum_{u,v}\mu_{U,V}(u,v)\log\mu_{U,V}(u,v)
 &=-\sum_{u,v}\mu_{U,V}(u,v)\log\mu_V(v)\\
 &\quad-\sum_{u,v}\mu_{U,V}(u,v)
                       \log\mu_{U\mid V=v}(u).
\end{align*}
Summing the first term over \(u\) gives \(\Ent(\mu_V)\).
Factoring \(\mu_V(v)\) out of the second for each \(v\) gives
the defined conditional entropy.
\end{proof}

Thus learning a pair can be separated into learning \(V\), then
learning what remains uncertain about \(U\) when \(V\) is known.
Independence is not assumed.

\section{Relative entropy: comparing two descriptions}\label{sec:KL}

\subsection{Definition, meaning, and support}

We will run the same tree under different probability models.
The quantity that compares these laws must involve two distributions.

\begin{definition}
For finite laws \(\mu,\nu\) on the same set, their \emph{relative
entropy}, or Kullback--Leibler divergence, is
\[
 \KL(\mu\Vert\nu)=\sum_{z:\mu(z)>0}\mu(z)\log\frac{\mu(z)}{\nu(z)}.
\]
If \(\nu(z)=0<\mu(z)\) for some \(z\), its value is \(+\infty\).
If every outcome of positive \(\mu\)-mass has positive \(\nu\)-mass,
we say that \(\mu\) is absolutely continuous with respect to \(\nu\).
\end{definition}

The logarithm measures how differently the two laws weight an outcome.
The average is under the \emph{first} law \(\mu\).
Individual logarithms can be negative, although the sum cannot be.
Entropy measures uncertainty within one law. Relative entropy measures
discrepancy between two laws, with an indicated order; it is not a
symmetric distance.
For example, \(\mu=(1,0)\), \(\nu=(1/2,1/2)\) give
\(\KL(\mu\Vert\nu)=1\) but \(\KL(\nu\Vert\mu)=+\infty\).
If \(\mu=\nu\), their divergence is zero.
If \(\mu\) is a point mass at \(z\), then
\begin{equation}\label{eq:pointKL}
 \KL(\mu\Vert\nu)=\log\frac1{\nu(z)}.
\end{equation}

\begin{lemma}[Nonnegativity]\label{lem:nonnegative}
\(\KL(\mu\Vert\nu)\ge0\).
\end{lemma}
\begin{proof}
Calculus gives \(-\ln s\ge1-s\) for \(s>0\): the difference has
derivative \(1-1/s\) and minimum zero at \(s=1\).
For finite divergence, apply this with \(s=\nu(z)/\mu(z)\):
\begin{align*}
 (\ln2)\KL(\mu\Vert\nu)
 &\ge\sum_{\mu(z)>0}(\mu(z)-\nu(z))\\
 &=1-\sum_{\mu(z)>0}\nu(z)\ge0.
\end{align*}
Infinite divergence is also nonnegative.
\end{proof}

Comparison with the uniform law on \(N\) outcomes gives
\(\KL(\mu\Vert\text{uniform})=\log N-\Ent(\mu)\).
Thus entropy on those outcomes is at most \(\log N\).
This statement is a consequence of the finite-sum proof, without
appealing to a separate information-theoretic theorem.

\subsection{The chain rule and the direction of comparison}

\begin{lemma}[Relative-entropy chain rule]\label{lem:KLchain}
For finite joint laws with \(\mu_{U,V}\) absolutely continuous with
respect to \(\nu_{U,V}\),
\begin{equation}\label{eq:KLchain}
 \KL(\mu_{U,V}\Vert\nu_{U,V})
 =\KL(\mu_V\Vert\nu_V)
 +\sum_{v:\mu_V(v)>0}\mu_V(v)
       \KL(\mu_{U\mid V=v}\Vert\nu_{U\mid V=v}).
\end{equation}
In particular, the marginal divergence is at most the joint divergence.
\end{lemma}
\begin{proof}
On positive \(\mu\)-mass pairs, factor the likelihood ratio:
\[
 \frac{\mu_{U,V}(u,v)}{\nu_{U,V}(u,v)}
 =\frac{\mu_V(v)}{\nu_V(v)}
  \frac{\mu_{U\mid V=v}(u)}{\nu_{U\mid V=v}(u)}.
\]
Take logarithms and average under \(\mu\). The first term sums to
the marginal divergence. In the second, average first under the
\(\mu\)-conditional law for each \(v\), then with weight \(\mu_V(v)\).
This is the displayed identity. Nonnegativity of each conditional
divergence proves the last assertion.
\end{proof}

We will use \(U\) for the hidden input and \(V\) for its transcript.
The identity separates discrepancy visible in the transcript from
discrepancy remaining in the conditional input law. If the joint
divergence is known and the conditional term has a lower bound,
subtracting that lower bound gives an \emph{upper} bound on transcript
divergence. We will write the identity again when its terms have
concrete meanings.

If \(V=f(U)\) is the same deterministic function under both laws,
the joint masses are the input masses attached to \((u,f(u))\).
Therefore
\(\KL(\mu_{U,V}\Vert\nu_{U,V})=\KL(\mu_U\Vert\nu_U)\).
The chain rule shows that applying the function cannot increase
divergence. Truthful execution of a fixed tree is one such function.

\section{Adding information along a decision tree}\label{sec:treeinformation}

\subsection{Why reach probabilities are the weights}

Fix a finite binary tree. A law \(\mu_Y\) on terminal transcripts
induces a reach probability \(\mu(v)\) at each internal node.
For \(\mu(v)>0\), let \(\pi_v\) be its conditional negative-answer
probability. Nodes of zero reach probability contribute zero.
Let \(T\) be the random number of tests.

\begin{lemma}[Tree decompositions]\label{lem:tree}
One has
\begin{align}
 \E_\mu T&=\sum_v\mu(v),\label{eq:tree-depth}\\
 \Ent(\mu_Y)&=\sum_v\mu(v)h(\pi_v).\label{eq:tree-entropy}
\end{align}
For another transcript law \(\nu_Y\), with conditional negative
probabilities \(\rho_v\), and with \(\mu_Y\) absolutely continuous
with respect to \(\nu_Y\),
\begin{equation}\label{eq:tree-divergence}
 \KL(\mu_Y\Vert\nu_Y)=\sum_v\mu(v)\kl(\pi_v\Vert\rho_v),
\end{equation}
where
\[
 \kl(u\Vert w)=u\log\frac uw+(1-u)\log\frac{1-u}{1-w}
\]
is the relative entropy of two binary laws, with the same
zero-probability and infinite-divergence conventions as above.
\end{lemma}
\begin{proof}
On each path, \(T\) is the sum of the indicators of visiting internal
nodes. Taking expectations proves \eqref{eq:tree-depth}.

The probability of a terminal transcript is the product of the
conditional probabilities of its edges. Expand the negative logarithm
of that product and average over terminal transcripts. A node \(v\)
appears only on paths that reach it, with total probability \(\mu(v)\).
Conditional on reaching it, its average contribution is
\(-\pi_v\log\pi_v-(1-\pi_v)\log(1-\pi_v)=h(\pi_v)\).
This proves \eqref{eq:tree-entropy}.
For \eqref{eq:tree-divergence}, expand the logarithm of the ratio of
the two products. The conditional average at a reached node is
\(\kl(\pi_v\Vert\rho_v)\). All sums are finite, so their order can
be changed without a convergence issue.
\end{proof}

The reach weight is indispensable. A test available only after a rare
history contributes only when that history occurs. An unweighted sum
over nodes would pretend that every branch is traversed in one run.

\begin{example}
Give the leaves of a binary tree depths \(1,2,2\) and probabilities
\(1/2,1/4,1/4\). The root is reached with probability one and splits
fairly. The other internal node is reached with probability \(1/2\)
and also splits fairly. Both expected depth and entropy are
\(1+(1/2)\cdot1=3/2\).
This illustrates a finite-tree identity; it does not assert that
arbitrary probability splits are realizable by pooled queries.
\end{example}

\subsection{Expected depth beyond the counting requirement}

In a correct tree on \(\binom{[n]}d\), choose the input uniformly.
Distinct valid inputs have distinct terminal transcripts, so the
transcript law is uniform on \(\binom nd\) outcomes. Its entropy is
\(H=\log\binom nd\).
Let \(P\) denote this truthful exact-count experiment, and put
\[
 L=\E_P T,\qquad \Gamma=L-H.
\]
The tree identities give
\begin{equation}\label{eq:Gamma}
 \Gamma=\sum_v P(v)\bigl(1-h(\pi_v)\bigr)\ge0,
 \qquad \pi_v=P(\text{answer }0\mid\text{reach }v).
\end{equation}
Thus \(\Gamma\) is both the expected test count in excess of \(H\)
and the visit-weighted sum of the amounts by which answers fall
short of one bit of entropy.

Averages satisfy \(L\le\max_D T(D)\), so a lower bound on \(L\)
is also a lower bound on that maximum. But a modest improvement over
\(H\) may still be below \(n-1\). The structural part of the proof
will be needed for the exact conclusion.

\subsection{A weighted comparison for one answer}

Suppose the negative probability is \(\pi\) in the actual experiment
and \(\rho\) in a comparison experiment, with \(\rho\) far from \(1/2\).
If \(\pi\) is also far from \(1/2\), the deficiency \(1-h(\pi)\)
is substantial. If \(\pi\) is near \(1/2\), it differs from \(\rho\),
producing relative entropy. The following lemma quantifies this
alternative.

\begin{lemma}[Weighted binary inequality]\label{lem:weighted}
For \(\omega>0\) and \(\pi,\rho\in[0,1]\),
\begin{equation}\label{eq:weighted}
 \omega\bigl(1-h(\pi)\bigr)+\kl(\pi\Vert\rho)
 \ge\frac{2\omega}{\omega+1}\frac{(\rho-1/2)^2}{\ln2}.
\end{equation}
\end{lemma}
\begin{proof}
For \(0<\rho<1\), let \(f(u)=\kl(u\Vert\rho)\). Differentiation gives
\[
 f(\rho)=f'(\rho)=0,\qquad
 f''(u)=\frac1{(\ln2)u(1-u)}\ge\frac4{\ln2}.
\]
Integrate twice from \(\rho\) to \(u\), reversing directions when
\(u<\rho\), to get \(f(u)\ge2(u-\rho)^2/\ln2\).
Apply this also with \(\rho=1/2\), noting
\(\kl(u\Vert1/2)=1-h(u)\). The left side of \eqref{eq:weighted}
is then at least
\[
 \frac2{\ln2}\left[\omega(\pi-1/2)^2+(\pi-\rho)^2\right].
\]
Complete the square:
\[
 \omega(\pi-1/2)^2+(\pi-\rho)^2
 =(\omega+1)\left(\pi-\frac{\omega/2+\rho}{\omega+1}\right)^2
  +\frac{\omega}{\omega+1}(\rho-1/2)^2.
\]
Discarding the square proves the result. Continuity handles finite
endpoint cases; an infinite relative entropy satisfies the inequality
as well.
\end{proof}

The weight \(\omega\) balances the two terms when summing over node
classes. It changes neither the algorithm nor a probability law.
We will use \(1/10\) for large pools and \(64\) for certain pairs.
Their role is to make the final coefficient calculation work, with
every numerical comparison checked exactly.

\section{Independent inputs and the cost of knowing the exact count}
\label{sec:product}

The real problem fixes the number of defectives. This section introduces
an independent input law for comparison and quantifies the cost of
conditioning it on the required count. At first we discuss input laws
only; executions on other counts will be defined in the next section.

\subsection{Two input laws on the same finite set}

Let \(X\subseteq[n]\) be a random subset, with indicators
\(X_i=\mathbf1\{i\in X\}\). Set \(p=d/n\) and \(q=1-p\).
There are two ways to choose \(X\):
\begin{itemize}
\item Under \(P_X\), choose uniformly among the \(d\)-element subsets.
Thus \(P_X(D)=1/\binom nd\) for \(|D|=d\), and zero otherwise.
\item Under \(Q_X\), choose the indicators independently, each equal
to one with probability \(p\). Thus
\[
 Q_X(D)=p^{|D|}q^{n-|D|}\qquad(D\subseteq[n]).
\]
\end{itemize}
The second is a Bernoulli product law. Its expected defective count
is \(np=d\), but its count is not fixed. Having the same mean as the
promised count does not make it the original input law.

Under \(P_X\), each item individually has defective probability \(p\),
but items are generally dependent. With exactly one defective among four,
\[
 P_X(X_2=1)=\frac14,\qquad
 P_X(X_2=1\mid X_1=1)=0.
\]
Under the independent law with \(p=1/4\), the latter probability
would still be \(1/4\).

All \(d\)-subsets have the same \(Q_X\)-mass \(p^dq^{n-d}\).
Consequently
\begin{equation}\label{eq:conditionrank}
 P_X=Q_X(\,\cdot\mid |X|=d).
\end{equation}
This is an identity, not an approximation.
For independent subsets,
\[
 \Ent(Q_X)
 =-\E_{Q_X}[|X|\log p+(n-|X|)\log q]=nh(p),
\]
because \(\E_{Q_X}|X|=np\).
Since \(\Ent(P_X)=H=\log\binom nd\), define
\begin{equation}\label{eq:Kdefinition}
 K=nh(p)-H.
\end{equation}
The next calculation explains why this particular difference also
has a divergence interpretation.

\begin{lemma}[Exact-count comparison]\label{lem:rank}
For the laws just defined,
\begin{equation}\label{eq:rankcost}
 \KL(P_X\Vert Q_X)=-\log Q_X(|X|=d)=K,\qquad
 0\le K\le\log(n+1).
\end{equation}
\end{lemma}
\begin{proof}
For every valid input \(D\),
\[
 \frac{P_X(D)}{Q_X(D)}
 =\frac1{\binom nd p^dq^{n-d}}
 =\frac1{Q_X(|X|=d)}.
\]
The logarithm of this ratio is constant on the support of \(P_X\).
Its average equals that same logarithm, namely
\(-H-d\log p-(n-d)\log q=nh(p)-H\).

For \(0\le k\le n\), put
\(b_k=Q_X(|X|=k)=\binom nk p^kq^{n-k}\).
For \(k<n\),
\[
 \frac{b_{k+1}}{b_k}=\frac{n-k}{k+1}\frac{d}{n-d}.
\]
This ratio decreases with \(k\). At \(k=d-1\) it equals
\((n-d+1)/(n-d)>1\); at \(k=d\) it equals \(d/(d+1)<1\).
Hence \(b_d\) is a largest one of \(b_0,\ldots,b_n\).
These \(n+1\) numbers sum to one, so \(1/(n+1)\le b_d\le1\).
Taking negative logarithms proves the bounds.
\end{proof}

\begin{example}
For \(n=4,d=1,p=1/4\),
\[
 Q_X(|X|=1)=4\left(\frac14\right)\left(\frac34\right)^3
 =\frac{27}{64}.
\]
Conditioning gives the uniform law on the four defective labels.
The comparison cost is \(K=\log(64/27)\).
The equality of costs across valid inputs is important; it does not
assert a small cost for every later conditioning event.
\end{example}

In particular, \(K\le\log(n+1)\) is one bound for the entire input-law
comparison. It is not a new bound that may be spent independently at
each node. The tree chain rule will keep the accumulated discrepancy
within this single total.

\subsection{Increasing and decreasing events}

Identify subsets with their \(0\)-\(1\) vectors. A function is
\emph{increasing} if changing any coordinate from zero to one cannot
decrease its value. An event is increasing if its indicator is
increasing. It is \emph{decreasing} if adding defectives cannot turn
a false event into a true one.

For fixed \(S\), the event \(X\cap S\ne\varnothing\) is increasing,
and \(X\cap S=\varnothing\) is decreasing.
An intersection of positive-test events is increasing too: adding
defectives preserves all the positive answers.

\begin{lemma}[Association for independent coordinates]\label{lem:association}
Under a product law on finitely many independent Bernoulli coordinates,
increasing real functions \(f,g\) satisfy \(\E(fg)\ge(\E f)(\E g)\).
If \(I\) is increasing with positive probability and \(J_0\) is
decreasing, then
\begin{equation}\label{eq:association}
 \Pr(J_0\mid I)\le\Pr(J_0).
\end{equation}
\end{lemma}
\begin{proof}
Define \(\operatorname{Cov}(f,g)=\E(fg)-(\E f)(\E g)\).
Induct on the number of coordinates. With no coordinates the
functions are constant and covariance is zero.
Let \(Z\) be the last coordinate. Conditional on \(Z\), the other
coordinates remain independent, so the induction hypothesis gives
\(\operatorname{Cov}(f,g\mid Z)\ge0\).
Put \(f_i=\E(f\mid Z=i)\), \(g_i=\E(g\mid Z=i)\).
The other-coordinate law is the same at \(i=0\) and \(1\), so
pointwise increase gives \(f_1\ge f_0\), \(g_1\ge g_0\).
If \(\Pr(Z=1)=u\), expansion gives
\[
 \operatorname{Cov}(f_Z,g_Z)=u(1-u)(f_1-f_0)(g_1-g_0)\ge0.
\]
Add and subtract \(\E[\E(f\mid Z)\E(g\mid Z)]\) to obtain
\[
 \operatorname{Cov}(f,g)
 =\E\operatorname{Cov}(f,g\mid Z)
  +\operatorname{Cov}(\E(f\mid Z),\E(g\mid Z)).
\]
Both terms are nonnegative. Constant coordinates may be omitted.

Finally, \(\mathbf1_I\) and \(-\mathbf1_{J_0}\) are increasing.
Their nonnegative covariance gives
\(\Pr(I\cap J_0)\le\Pr(I)\Pr(J_0)\).
Divide by \(\Pr(I)>0\).
\end{proof}

Coordinates need not remain independent after a positive answer.
The lemma says instead that increasing constraints cannot raise the
probability of a decreasing event. That inequality is what we need.

\section{Preparing one tree and defining three experiments}\label{sec:experiments}

Fix a correct deterministic tree for \(\binom{[n]}d\). The tree is
not random. We will run it under three laws. First we make its
behavior unambiguous even on histories impossible under the promise.

\subsection{Completing missing branches}

If an answer branch was omitted because no valid input can take it,
attach an immediate leaf with an arbitrary output. This changes no
valid input path. It only provides an endpoint for a comparison
experiment that takes a formerly omitted edge.
The completed tree is finite. It need not identify inputs outside
the exact-\(d\) promise; we only need all experiments to be defined.

\subsection{Removing labels already declared negative}

At a node \(v\), let \(Z_v\) be the union of all original pools
receiving negative answers earlier on its prefix transcript.
Replace its original query \(S_v\) by
\begin{equation}\label{eq:normalization}
 \widehat S_v=S_v\setminus Z_v.
\end{equation}
This is the \emph{normalized pool}. The set \(Z_v\) is determined by
the edge labels on the path, before any probability law is chosen.

\begin{lemma}\label{lem:normalization}
Normalization preserves every truthful execution on every input subset.
\end{lemma}
\begin{proof}
At the root \(Z_v=\varnothing\). Suppose the original and normalized
executions have reached the same node after the same truthful answers.
Every label in a previous negative pool is nondefective, so
\(X\cap Z_v=\varnothing\). Hence
\[
 X\cap S_v=\varnothing
 \quad\Longleftrightarrow\quad X\cap(S_v\setminus Z_v)=\varnothing.
\]
The answers agree. After a negative edge, updating the union by the
normalized or original pool gives the same result because
\(Z_v\cup(S_v\setminus Z_v)=Z_v\cup S_v\).
Induction proves agreement through termination.
\end{proof}

This proof works for all input subsets, even those of size different
from \(d\). Empty normalized pools may be retained.
From now on, pool size and label appearance refer to normalized
pools. A node is \emph{small} if \(|\widehat S_v|\le2\) and
\emph{large} if \(|\widehat S_v|\ge3\).
This classification is fixed by the node and its history.

\subsection{The truthful experiments}

Under \(P_{X,Y}\), draw \(X\) from \(P_X\) and answer all normalized
queries truthfully. This is the distribution of valid executions
used in the lower bound.
Under \(Q_{X,Y}\), draw \(X\) from \(Q_X\) and also answer truthfully.
The completed tree defines the execution even when \(|X|\ne d\).
Under these two laws \(Y\) is the same deterministic function of \(X\).
The exact-count comparison and the chain rule therefore give
\begin{equation}\label{eq:truthfulcomparison}
 \KL(P_{X,Y}\Vert Q_{X,Y})=K,\qquad \KL(P_Y\Vert Q_Y)\le K.
\end{equation}
This controls accumulated discrepancy over the entire transcript.

\subsection{The artificial reference experiment}

The third joint law \(R_{X,Y}\) is defined by:
\begin{enumerate}
\item Draw \(X\) from the independent law \(Q_X\).
\item At a small node, answer with \(\OR_X(\widehat S_v)\).
\item At a large node, ignore its OR value and use a new fair bit,
independent of the input and all preceding bits.
\item Follow that edge and repeat until a leaf is reached.
\end{enumerate}
This is a finite experiment. It could be implemented by sampling a
fair bit for every large node in advance and using it if visited.
The law \(R_{X,Y}\) records only input and transcript; unused bits
are summed out.

The next query remains adaptive, selected by the fixed tree from the
answers already supplied. Normalization still follows
\eqref{eq:normalization} syntactically.
An artificial negative on a large pool removes its labels from later
queries, but does \emph{not} establish that their input indicators
are zero under \(R\).
A small answer is truthful for the normalized pool, which need not
have the original pool's OR value under an artificial history.

Consequently a reference transcript may be inconsistent with its
sampled input. This is legitimate: \(R\) is a comparison measure,
not a proposed identification algorithm. The real algorithm under
\(P\) still receives only truthful answers and no additional oracle.

\begin{center}
\begin{tabularx}{\linewidth}{lYYY}
\toprule
Joint law&Input law&Answers in the normalized tree&Purpose\\
\midrule
\(P_{X,Y}\)&Uniform on \(|X|=d\)&All truthful&
Average over valid inputs\\
\(Q_{X,Y}\)&Independent Bernoulli-\(p\)&All truthful&
Control large-pool probabilities\\
\(R_{X,Y}\)&Independent Bernoulli-\(p\)&Truthful for size at most two;
fair independent bits for larger pools&
Control fresh pairs and untouched labels\\
\bottomrule
\end{tabularx}
\end{center}
The \(Q\) and \(R\) input marginals agree, but their transcript and
joint laws generally differ. All averaged counts in the main
argument will use \(P\), not \(Q\) or \(R\).

\subsection{A fixed-history formula}\label{subsec:history}

Fix a prefix \(\tau\) leading to a node. Let \(a(\tau)\) count its
preceding large nodes. Let \(E_\tau\) be the event that \(X\) gives
the prescribed answers at every preceding \emph{small normalized}
pool. These pools are fixed once \(\tau\) is fixed.
For every \(D\subseteq[n]\),
\begin{equation}\label{eq:historyformula}
 R(X=D,\ \text{reach }\tau)
 =Q_X(D)\,2^{-a(\tau)}\,\mathbf1_{E_\tau}(D).
\end{equation}
Indeed, each prescribed large answer has probability \(1/2\);
each prescribed small answer has probability one or zero according
to whether \(D\) satisfies it. Multiply these factors along the
fixed path, together with the input mass.
The same formula holds at a terminal transcript, with reaching it
meaning termination there.

Sum over \(D\), and then divide when the reach probability is positive:
\[
 R(\text{reach }\tau)=2^{-a(\tau)}Q_X(E_\tau),\qquad
 R(X=D\mid\text{reach }\tau)=Q_X(D\mid E_\tau).
\]
In particular,
\begin{equation}\label{eq:posterior}
 R_{X\mid\text{reach }\tau}=Q_X(\,\cdot\mid E_\tau).
\end{equation}

Let \(V_\tau\) be the union of the preceding small pools.
The event \(E_\tau\) depends only on coordinates in \(V_\tau\).
For a configuration \(z\) on its complement, independence gives
\[
 Q_X((X_i)_{i\notin V_\tau}=z,\ E_\tau)
 =Q_X((X_i)_{i\notin V_\tau}=z)\,Q_X(E_\tau).
\]
After division, those coordinates have their original product law
under the conditional reference experiment. Thus labels absent from
all preceding small pools remain mutually independent Bernoulli-\(p\),
even if they occurred in large pools. No independence within
\(V_\tau\) is asserted.

\begin{example}[An artificial negative]
Take six labels with promised count \(d=2\), so that \(p=1/3\).
Suppose the first query is \(\{1,2,3\}\), and after
answer zero the original next query is \(\{2,4,5\}\). Its normalized
version is \(\{4,5\}\). These can be initial queries of a legal
strategy that finishes by individual testing on each branch.

Under \(R\), the first answer is fair and can be zero even for
the sampled input \(D=\{1,4\}\). The second normalized query then
answers one truthfully. On the fixed prefix \(01\), the only input
constraint is \(X\cap\{4,5\}\ne\varnothing\), and
\[
 R(X=D,\text{reach }01)
 =Q_X(D)\,\frac12\,\mathbf1\{D\cap\{4,5\}\ne\varnothing\}.
\]
Coordinates \(1,2,3,6\) remain independent Bernoulli-\(p\) conditional
on this prefix. The artificial zero has not fixed \(1,2,3\) to be
nondefective. Coordinates \(4,5\), however, are constrained by a positive OR
answer and need not be independent. This illustrates the experiment,
not an application of the forthcoming density-\(1/4\) inequality.
\end{example}

\section{The whole-tree information bound near density one quarter}
\label{sec:whole}

We now apply the tools to a fixed correct tree. Retain \(p=d/n\),
\(q=1-p\), and suppose throughout this section that
\begin{equation}\label{eq:localinterval}
 \frac{249}{1000}\le p\le\frac{251}{1000}.
\end{equation}
Near \(p=1/4\), three independent items are all good
(that is, nondefective) with probability
near \((3/4)^3<1/2\), whereas two are all good with probability near
\((3/4)^2>1/2\). This explains the division into large pools and
pairs. The actual argument must make such bounds valid after
adaptive histories, not only before any tests.

\begin{theorem}[The whole-tree inequality]\label{thm:whole}
Suppose \eqref{eq:localinterval} holds. For any finite correct
deterministic adaptive tree, set
\[
 L=\E_P T,\qquad H=\log\binom nd,\qquad
 K=nh(p)-H,\qquad \Gamma=L-H.
\]
Then
\begin{equation}\label{eq:WT}
 \boxed{6615\Gamma\ge n-H-67229K.}
\end{equation}
The same inequality holds with \(M(d,n)-H\) in place of \(\Gamma\).
\end{theorem}

When the right side is positive, this forces a positive expected
excess over \(H\), and is stronger than the information bound alone.
Its form will matter in the final proof: \(K\le\log(n+1)\) is a
logarithmic error, whose size we will control explicitly.
We prove the theorem in the following stages.

Here is the accounting plan. We will introduce three expected counts:
\(A\), the number of large queries; \(B\), the number of pairs whose
two labels have not appeared in a preceding small query; and \(W\),
the number of labels never appearing in a small query. Counting
first appearances separately on each path will give
\[
 n-L\le B+W-A.
\]
Thus we need upper bounds for \(A\) and for \(B+W\). Comparing \(P\)
with the truthful product law \(Q\) will control \(A\). Comparing
\(P\) with the artificial law \(R\) will control \(B+W\), while
retaining a useful subtraction for the large nodes. The separate
entropy deficiencies will then combine within \(\Gamma\).
This paragraph describes the destination; each inequality is proved below.

\subsection{Separate the entropy contributions}

Recall
\(\Gamma=\sum_vP(v)(1-h(\pi_v))\), with \(\pi_v\) the actual
\(P\)-conditional negative probability.
In both transcript comparisons, the first law in the divergence is
\(P_Y\). The tree-divergence formula therefore uses the contributions
\[
 P(v)\kl(\pi_v\Vert\rho_v^Q)
 \quad\hbox{and}\quad
 P(v)\kl(\pi_v\Vert\rho_v^R),
\]
respectively. The reach weight is always \(P(v)\), not \(Q(v)\)
or \(R(v)\). Only the second binary parameter changes with the
comparison experiment.
Define
\[
 A=\sum_{v\ \mathrm{large}}P(v),\qquad
 \Glg=\sum_{v\ \mathrm{large}}P(v)(1-h(\pi_v)).
\]
The first is the expected number of large queries, while the second
is their contribution to entropy deficiency. Both use actual-law
weights. We will keep these contributions separate from those of a
second query class, rather than giving every class the full sum
\(\Gamma\).

\subsection{Large pools under the truthful product law}

Every \(P\)-reachable node is \(Q\)-reachable, since some valid
input reaches it and every input subset has positive \(Q_X\)-mass.
Fix such a node \(v\), and work under the truthful experiment \(Q\).
Reaching it imposes two kinds of conditions. First, every coordinate
in \(Z_v\), the union of earlier negative pools, must be zero.
Conditional on these fixed zeros, the remaining coordinates still
have their independent Bernoulli-\(p\) law. Second, each earlier
positive answer requires its pool to contain a defective. After
deleting coordinates in \(Z_v\), the intersection \(I_v\) of those
positive requirements is an increasing event on the remaining
coordinates.

Write \(F_v=\{X_i=0\text{ for every }i\in Z_v\}\), and let
\(J_v=\{X\cap\widehat S_v=\varnothing\}\).
The latter event is decreasing and uses only coordinates outside
\(Z_v\). Reaching \(v\) is exactly the event \(F_v\cap I_v\);
it has positive \(Q\)-probability. Apply the association inequality
under the product law conditional on \(F_v\):
\begin{equation}\label{eq:largeprob}
 \begin{aligned}
 \rho_v^Q
 &:=Q(\text{answer }0\mid\text{reach }v)
   =Q(J_v\mid F_v,I_v)\\
 &\le Q(J_v\mid F_v)=q^{|\widehat S_v|}.
 \end{aligned}
\end{equation}
Overlap among the positive pools, or between them and the current
pool, is allowed. The remaining coordinates need not be independent
after conditioning on \(I_v\); association is what justifies the
inequality without that additional independence.

At a large node,
\[
 \rho_v^Q\le(751/1000)^3<\frac{17}{40},\qquad
 \frac12-\rho_v^Q>\frac3{40}.
\]
The cubic inequality is certified by
\(751^3=423564751<425000000\).
Also \(\ln2<7/10\), because
\[
 e^{7/10}>1+\frac7{10}+\frac{49}{200}+\frac{343}{6000}
 =\frac{12013}{6000}>2.
\]
The strict lower bound on the exponential follows, for example,
by successive integration of \(e^x>1\) for \(x>0\).

Use the weighted inequality with \(\omega=1/10\).
Its right side is greater than
\[
 \frac2{11}\frac{(3/40)^2}{7/10}
 =\frac9{6160}>\frac1{685},
 \qquad 9\cdot685=6165>6160.
\]
At every large node with \(P(v)>0\), multiplication by \(P(v)\) gives
\[
 \frac{P(v)}{685}
 <\frac{P(v)}{10}(1-h(\pi_v))
   +P(v)\kl(\pi_v\Vert\rho_v^Q).
\]
Sum over large nodes. The first terms total \(\Glg/10\);
the second terms form part of the nonnegative tree-divergence sum
\(\KL(P_Y\Vert Q_Y)\le K\). Nodes with \(P(v)=0\) add zero.
Weakening the resulting strict bound if necessary, we obtain
\begin{equation}\label{eq:Alarge}
 \frac A{685}\le\frac{\Glg}{10}+K,\qquad
 \boxed{A\le\frac{137}{2}\Glg+685K.}
\end{equation}
A large expected number of large queries thus requires entropy
deficiency on those nodes or comparison cost. The latter is limited
by the single whole-tree quantity \(K\).

\subsection{Fresh pairs under the reference law}

A \emph{fresh pair} is a normalized query with two labels neither
of which has appeared in a preceding small normalized pool.
They may have appeared in large pools. Define
\[
 B=\sum_{v\ \mathrm{fresh}}P(v),\qquad
 \Gfr=\sum_{v\ \mathrm{fresh}}P(v)(1-h(\pi_v)).
\]
Thus \(B\) is an expected number of queries and \(\Gfr\) their
entropy contribution. Fresh pairs are small, so their nodes are
disjoint from large nodes and
\begin{equation}\label{eq:entropysplit}
 \Glg+\Gfr\le\Gamma.
\end{equation}

Every \(P\)-reachable history is also \(R\)-reachable. To see this,
choose a valid input reaching it truthfully. This input has positive
\(Q_X\)-mass, it satisfies every small-answer constraint, and the
finite prescribed sequence of fair large answers has positive
probability. Formula~\eqref{eq:historyformula} therefore gives positive
reference mass. The same argument at terminal transcripts proves
that \(P_Y\) is absolutely continuous with respect to \(R_Y\).
It also shows that every pair of positive \(P_{X,Y}\)-mass has
positive \(R_{X,Y}\)-mass. All subsequent reference conditional
probabilities and divergences are consequently well defined.

The reference posterior \eqref{eq:posterior} leaves the two fresh
coordinates independent Bernoulli-\(p\). Consequently
\[
 \rho_v^R=R(\text{answer }0\mid\text{reach }v)=q^2
 \ge(749/1000)^2>\frac{14}{25},\qquad
 \rho_v^R-\frac12>\frac3{50}.
\]
The exact square comparison is \(749^2=561001>560000\).
Earlier large positive answers in the truthful experiment might have
constrained these labels. In \(R\), those answers impose no input
constraint on a fixed history; this is why the reference process is
useful here.

\subsection{The joint divergence has an exact value}

For a valid input \(D\), let \(y(D)\) be its truthful terminal
transcript, \(\ell(D)\) its terminal leaf, and \(a(D)\) the number
of large queries along that path.
The actual joint law gives \((D,y(D))\) mass \(1/\binom nd=2^{-H}\).
In the reference experiment every small answer on this truthful
path agrees with \(D\), and every prescribed large answer has
probability \(1/2\). Thus
\begin{equation}\label{eq:Rpoint}
 R_{X,Y}(D,y(D))=Q_X(D)\,2^{-a(D)}
 =2^{-nh(p)-a(D)}.
\end{equation}
The logarithmic ratio of actual to reference joint mass there is
\[
 \log\frac{2^{-H}}{2^{-nh(p)-a(D)}}=K+a(D).
\]
Average under \(P\), for which \(\E_P a(X)=A\), to obtain
\begin{equation}\label{eq:jointvalue}
 \boxed{\KL(P_{X,Y}\Vert R_{X,Y})=K+A.}
\end{equation}
This is a statement about input and transcript together. It is not
yet a value for transcript divergence alone.

\subsection{Labels untouched by small pools}

On the truthful path of \(D\), let \(V(D)\) be the union of all
small normalized pools. Define
\[
 w_1(D)=|D\setminus V(D)|,\quad
 w_0(D)=|([n]\setminus D)\setminus V(D)|,\quad
 w(D)=w_1(D)+w_0(D).
\]
These count defective, nondefective, and total labels never seen in
a small pool. They are quantities on one path, not yet averages.
They may include labels used in large pools. Put \(W=\E_P w(X)\).

Consider the following particular class of reference inputs:
\[
 \mathcal A_D=\{D'\subseteq[n]:D'\cap V(D)=D\cap V(D)\}.
\]
Each agrees coordinate by coordinate with \(D\) on every small pool
of the fixed transcript \(y(D)\), so supplies all the prescribed
small answers. The prescribed large answers depend only on fair
reference bits and together have probability \(2^{-a(D)}\).
Thus the total reference mass reaching \(y(D)\) from this class is
\[
 2^{-a(D)}Q_X(\mathcal A_D)
 =2^{-a(D)}p^{d-w_1(D)}q^{n-d-w_0(D)}.
\]
Here we fixed \(d-w_1(D)\) defective and \(n-d-w_0(D)\) nondefective
coordinates inside \(V(D)\). Each outside coordinate is unrestricted:
summing its two possibilities contributes \(p+q=1\). There may be
additional reference inputs producing the prescribed small OR
answers without agreeing coordinate by coordinate on \(V(D)\).
Ignoring those inputs gives a lower bound:
\begin{align}
 R_Y(y(D))
 &\ge2^{-a(D)}p^{d-w_1(D)}q^{n-d-w_0(D)}\nonumber\\
 &=2^{-a(D)}Q_X(D)\,p^{-w_1(D)}q^{-w_0(D)}.
 \label{eq:leafbound}
\end{align}
Inputs in \(\mathcal A_D\) can have different total defective counts.
This is permitted in the reference experiment; the real algorithm's
exact-count promise has not been enlarged.

\subsection{One bound shared by transcript and conditional divergence}

Write
\[
 \mathcal D_R=\KL(P_Y\Vert R_Y),\qquad
 J=\E_P\left[w_1(X)\log\frac1p+w_0(X)\log\frac1q\right].
\]
Here \(\mathcal D_R\) concerns transcripts alone. The quantity \(J\)
is an actual-law average of information associated with untouched
labels. The joint chain rule is the exact identity
\begin{equation}\label{eq:mainchain}
 K+A=\mathcal D_R+
 \sum_{y:P_Y(y)>0}P_Y(y)
       \KL(P_{X\mid Y=y}\Vert R_{X\mid Y=y}).
\end{equation}

At \(y(D)\), correctness leaves exactly one valid input, \(D\).
Thus the actual conditional input law is the point mass at \(D\).
The reference conditional probability is
\[
 R(X=D\mid Y=y(D))
 =\frac{R_{X,Y}(D,y(D))}{R_Y(y(D))}
 \le p^{w_1(D)}q^{w_0(D)}.
\]
The inequality follows from \eqref{eq:Rpoint} and \eqref{eq:leafbound}:
the numerator is known exactly and the denominator has a lower
bound, producing an upper bound on the quotient. By the point-mass
formula~\eqref{eq:pointKL}, the conditional divergence is at least
\[
 w_1(D)\log\frac1p+w_0(D)\log\frac1q.
\]
Substitute its average into \eqref{eq:mainchain}:
\begin{equation}\label{eq:shared}
 0\le\mathcal D_R,\qquad \boxed{\mathcal D_R+J\le K+A.}
\end{equation}

The direction has a concrete meaning. The full joint divergence is
\(K+A\). At least \(J\) of it remains in the conditional input
comparison after the transcript is observed. Hence at most
\(K+A-J\) is available for transcript divergence. The lower bound
on the conditional term produces an upper bound on the marginal
term, and this will let us bound two counts together.

The direction can also be checked directly from leaf probabilities:
\[
 \log\frac{P_Y(y(D))}{R_Y(y(D))}
 \le K+a(D)-w_1(D)\log\frac1p-w_0(D)\log\frac1q.
\]
Its \(P\)-average is exactly \(\mathcal D_R\le K+A-J\).

We need a simple lower bound on the information per untouched label.
In our interval \(p\le q\le751/1000<151/200\).
The exact inequalities
\[
 (151/200)^2<4/7,\qquad
 (151/200)^5<(4/7)^2(151/200)=302/1225<1/4
\]
show that \(p^5,q^5<2^{-2}\).
For the first rational comparison, \(7\cdot22801=159607<160000\);
for the last, \(4\cdot302=1208<1225\).
Taking logarithms gives \(\log(1/p)>2/5\) and \(\log(1/q)>2/5\).
Therefore
\begin{equation}\label{eq:J}
 J\ge\frac25 W.
\end{equation}

\subsection{Retain the contribution of fair artificial answers}

At a large node the conditional reference answer is fair. Its
contribution to the transcript divergence is
\(\kl(\pi_v\Vert1/2)=1-h(\pi_v)\).
The total contribution of large nodes to \(\mathcal D_R\) is
therefore exactly \(\Glg\).
All other node divergences are nonnegative. In particular the sum
over fresh-pair nodes is at most \(\mathcal D_R-\Glg\).

At a fresh pair use the weighted inequality with \(\omega=64\).
Its right side is greater than
\[
 \frac{128}{65}\frac{(3/50)^2}{7/10}
 =\frac{576}{56875}>\frac1{99},
 \qquad 576\cdot99=57024>56875.
\]
At a fresh-pair node with \(P(v)>0\), we consequently have
\[
 \frac{P(v)}{99}
 <64P(v)(1-h(\pi_v))
   +P(v)\kl(\pi_v\Vert\rho_v^R).
\]
Sum with the actual-law reach weights. The first terms total
\(64\Gfr\). The divergence terms from all large nodes total exactly
\(\Glg\), and every other node contributes nonnegatively.
The fresh-pair divergence sum is therefore at most
\(\mathcal D_R-\Glg\), yielding
\begin{equation}\label{eq:freshsum}
 \frac B{99}\le64\Gfr+\mathcal D_R-\Glg.
\end{equation}
Combine this with \eqref{eq:shared}, keeping its conditional term:
\[
 \frac B{99}+J\le64\Gfr+K+A-\Glg.
\]
Multiply by \(99\). Since \(99J\ge(198/5)W\ge W\), we get
\begin{equation}\label{eq:coupled}
 \boxed{B+W\le6336\Gfr+99K+99A-99\Glg.}
\end{equation}
The estimate controls fresh pairs and untouched labels jointly.
The subtraction of \(\Glg\) is useful: the large nodes have already
used that exact portion of the transcript divergence, leaving only
the remainder for fresh pairs.

\subsection{Count first appearances on a single path}

Fix a valid \(D\). Let \(t(D)\) be its number of tests and \(b(D)\)
its number of fresh pairs. Its number of small queries is
\(t(D)-a(D)\).
Count a label when it first appears in a small pool:
\begin{itemize}
\item an empty pool adds none;
\item a singleton adds at most one;
\item a pair with a previously seen label adds at most one;
\item a fresh pair adds two, one more than that general allowance.
\end{itemize}
Consequently
\[
 |V(D)|\le t(D)-a(D)+b(D).
\]
This is a pathwise combinatorial inequality, valid with arbitrary
overlaps. Since \(|V(D)|=n-w(D)\), average under \(P\) to obtain
\begin{equation}\label{eq:firstappear}
 n-W\le L-A+B,\qquad n-L\le B+W-A.
\end{equation}
The probability law is the same for all the expectations on this line.

\subsection{Finish the coefficient calculation}

Insert \eqref{eq:coupled} into \eqref{eq:firstappear}, and then use
\eqref{eq:Alarge}:
\begin{align*}
 n-L
 &\le6336\Gfr+99K+99A-99\Glg-A\\
 &=6336\Gfr+99K+98A-99\Glg\\
 &\le6336\Gfr+
       \left(98\cdot\frac{137}{2}-99\right)\Glg
       +(99+98\cdot685)K\\
 &=6336\Gfr+6614\Glg+67229K.
\end{align*}
The arithmetic is \(98\cdot137/2=6713\) and
\(98\cdot685=67130\).
Because \(6336\le6614\) and \(\Gfr+\Glg\le\Gamma\),
\[
 n-L\le6614\Gamma+67229K.
\]
Now \(L=H+\Gamma\), so \(n-H=(n-L)+\Gamma\) gives
\[
 n-H\le6615\Gamma+67229K,
\]
which is \eqref{eq:WT}. An optimal worst-case tree has
\(L\le M(d,n)\). Replacing \(L-H\) by the larger \(M(d,n)-H\)
preserves the lower-bound inequality and proves the last assertion
of Theorem~\ref{thm:whole}.\quad\hfill\(\square\)

All query types have been included: large pools through the product
comparison, fresh pairs through the reference posterior, and other
small pools through nonnegative omitted terms and the first-appearance
count. The finite-tree identities allow one-sided feasible answers,
empty pools, and unequal leaf depths. The artificial process has
served only as a comparison law.

\section{A structural relation between different testing problems}
\label{sec:structural}

The information bound applies near density \(1/4\). To use it for
Theorem~\ref{thm:main}, we need a relation that moves between complete
families of different sizes. This section is entirely deterministic.
Its simulations use actual pooled-OR queries; no independent input law
or artificial answers are involved.

\subsection{Complexity of a candidate family}

A \emph{candidate family} \(\mathcal F\) is a nonempty finite set of
possible defective subsets of a specified ground set. The algorithm
is told the family, and must identify its member using legal OR tests.
A legal query is still any subset \(S\) of that ground set, with
answer one on candidate \(D\) exactly when \(S\cap D\ne\varnothing\).
Knowing the family does not mean knowing which member is actual,
and does not permit arbitrary binary questions about its members.
Let \(\Comp(\mathcal F)\) denote the minimum worst-case depth.
Thus
\[
 M(d,n)=\Comp\left(\binom{[n]}d\right).
\]
This notation lets us discuss the families left after tests, which
need not be complete families of equal-size subsets.

\begin{lemma}[Elementary reductions]\label{lem:family}
The following operations have the stated effects:
\begin{enumerate}
\item Restriction to a nonempty subfamily cannot increase complexity.
\item Relabeling items preserves complexity.
\item Removing labels known to be nondefective preserves complexity.
\item Adjoining the same fixed defective labels to every candidate
preserves complexity.
\item A tree may be reduced so that both children of every internal
node have nonempty candidate families, without increasing its depth.
\end{enumerate}
\end{lemma}
\begin{proof}
A tree correct for a family is also correct for each subfamily;
impossible paths may be removed. This proves the first assertion
and fixes its direction.
For relabeling, apply the same bijection to every candidate and query.
All intersection answers are preserved, and the inverse bijection
gives the reverse comparison.

A known nondefective label can be deleted from any pool without
changing an answer. Conversely an algorithm that ignores that label
still works when it is present.
If every candidate contains a fixed defective label, any query
containing it is constantly positive. Suppress such a query, following
its positive branch without a test. Other queries ask only about
the remaining unknown part. In the other direction, a tree for that
unknown part also identifies the family with the fixed labels
adjoined. This proves equality of complexities.

Finally, if one child of an internal node has no possible candidates,
the answer is forced on every remaining input. Suppress that query
and continue into the other child. Repeating this operation terminates
in a finite tree. Leaves remain correct and path lengths do not increase.
\end{proof}

These reductions use information already contained in the specified
candidate family. They do not give the algorithm extra knowledge about
which member is the actual input.
In a reduced tree an internal node cannot have only one surviving
candidate, since then no query could have two nonempty children.

\subsection{The endpoint problems}

\begin{lemma}\label{lem:endpoints}
\[
 M(N-1,N)=N-1\quad(N\ge2),\qquad M(1,2)=1,\qquad M(1,3)=2.
\]
\end{lemma}
\begin{proof}
When exactly one item is good, every pool of size at least two is
positive. Empty queries are constant too. A singleton query can
ask whether a particular label is the one good item.
For the lower bound, follow a feasible adversarial branch. While at
least two possible locations of the nondefective item remain,
answer a singleton test positively whenever its label is still a
candidate location. This declares that label defective, eliminates
exactly one possible location, and leaves at least one candidate.
Pools of size at least two, and singletons already eliminated as
possible locations, are forced positive and eliminate nothing further.
An empty pool is forced negative and likewise eliminates nothing.
Starting with \(N\) possible locations, this feasible branch therefore
requires at least \(N-1\) tests before only one location remains.
Individual testing achieves that number.

With one defective among two or three items, the numbers of
candidates are respectively two and three. A binary tree needs at
least one or two tests by leaf counting. Individual testing achieves
those bounds.
\end{proof}

\subsection{Two distinguished labels}

The next identity is the key structural step. It explains the
complexity of a family in which we know that at least one of two
specified labels is defective.

Let \(U\) contain \(m\) ordinary labels, and let \(\xi,\eta\) be two
additional distinct labels. Fix \(1\le k<m\). Define
\[
 \mathcal F=\binom Uk,\qquad \mathcal G=\binom U{k-1}.
\]
For this subsection only, \(\xi D\) means \(\{\xi\}\cup D\), and
\(\xi\eta E\) means \(\{\xi,\eta\}\cup E\).
Consider
\begin{align*}
 \mathcal J
 &=\{\xi D:D\in\mathcal F\}
   \cup\{\eta D:D\in\mathcal F\}
   \cup\{\xi\eta E:E\in\mathcal G\},\\
 \mathcal C
 &=\mathcal F\cup\{\eta E:E\in\mathcal G\}
   =\binom{U\cup\{\eta\}}k.
\end{align*}
Every member of \(\mathcal J\) has exactly \(k+1\) defectives and
meets \(\{\xi,\eta\}\). Its three types are:
\begin{center}
\begin{tabular}{lll}
\toprule
Distinguished defectives&Ordinary defectives&Candidate form\\
\midrule
Only \(\xi\)&\(k\) labels from \(U\)&\(\xi D\)\\
Only \(\eta\)&\(k\) labels from \(U\)&\(\eta D\)\\
Both&\(k-1\) labels from \(U\)&\(\xi\eta E\)\\
\bottomrule
\end{tabular}
\end{center}
The family \(\mathcal C\) is an ordinary complete problem: \(k\)
defectives among \(m+1\) labels.

\begin{lemma}[Two-label identity]\label{lem:two}
\[
 \Comp(\mathcal J)=1+\Comp(\mathcal C)=1+M(k,m+1).
\]
\end{lemma}

\paragraph{The easy direction.}
Test \(\{\xi\}\). On a positive answer the remaining family is
\(\{\xi C:C\in\mathcal C\}\), whose complexity is
\(\Comp(\mathcal C)\) after deleting the fixed defective \(\xi\).
On a negative answer the family is
\(\{\eta D:D\in\mathcal F\}\), with complexity
\(\Comp(\mathcal F)\le\Comp(\mathcal C)\) by restriction.
Thus \(\Comp(\mathcal J)\le1+\Comp(\mathcal C)\).

The reverse direction must work even if a proposed tree never begins
by testing a distinguished label. We will convert any reduced tree
for \(\mathcal J\), of height \(s\), into a legal tree for
\(\mathcal C\) of height at most \(s-1\).

\paragraph{What can ordinary queries distinguish?}
A query contained in \(U\) will be called ordinary.
After a prefix made only of ordinary queries, let
\(\mathcal A\subseteq\mathcal F\) and
\(\mathcal B\subseteq\mathcal G\) be the ordinary subsets consistent
with those answers. The surviving family in \(\mathcal J\) is
\begin{equation}\label{eq:ordinaryfamily}
 \{\xi D:D\in\mathcal A\}
 \cup\{\eta D:D\in\mathcal A\}
 \cup\{\xi\eta E:E\in\mathcal B\}.
\end{equation}
In particular, \(\xi D\) and \(\eta D\) cannot be distinguished by
ordinary queries: both meet an ordinary pool exactly when \(D\) does.

There is also a useful relation between the two subset sizes.
If \(E,E'\) are distinct members of \(\mathcal B\), choose
\(z\in E'\setminus E\). Then \(E\cup\{z\}\) belongs to
\(\mathcal A\). To verify this, every earlier positive pool meets
\(E\), and therefore also meets \(E\cup\{z\}\).
Every earlier negative pool avoids \(E'\), hence avoids \(z\);
it also avoids \(E\). So all negative answers are preserved.
The enlarged set has size \(k\), as required. We have proved
\begin{equation}\label{eq:extension}
 \mathcal A=\varnothing,\quad \mathcal B\ne\varnothing
 \quad\Longrightarrow\quad |\mathcal B|=1.
\end{equation}

\paragraph{A deepest ordinary-only path is impossible.}
Suppose a reduced \(\mathcal J\)-tree of height \(s\) has a leaf
at depth \(s\) identifying \(\xi\eta E\), and all queries on its
path are ordinary. Consider its parent, at depth \(s-1\).
Its two nonempty children must both be leaves, so at most two
candidates can survive at the parent.
If its \(\mathcal A\) is nonempty, any \(D\in\mathcal A\) supplies
the distinct candidates \(\xi D,\eta D,\xi\eta E\), giving at
least three. If \(\mathcal A\) is empty, \eqref{eq:extension}
makes the surviving family a singleton, contrary to an internal
node of a reduced tree. Both possibilities are impossible.

Therefore an ordinary-only leaf of the form \(\xi\eta E\) has
depth at most \(s-1\). This observation will handle paths where
there is no special query to skip.

\paragraph{The simulation when a distinguished query occurs.}
The idea is to assign each possible input \(C\) of the new problem
an injective representative \(\phi(C)\) in the old problem.
Once this assignment is selected, the simulation maintains the
following invariant: every currently possible \(C\) has its
representative among the surviving candidates of the old tree,
and every query actually asked has the same answer on \(C\) as the
old query has on \(\phi(C)\). An old query may be skipped only when
its answer is known for every current representative. The new
transcript then follows a well-defined path through the old tree,
without knowing which input is actual. We will choose one of two
assignments at the first distinguished query; both preserve every
earlier ordinary answer.

Run a new algorithm on an unknown \(C\in\mathcal C\). Until the first
query containing a distinguished label, use the same ordinary queries
as the \(\mathcal J\)-tree. The possible inputs in \(\mathcal C\)
at that point are \(\mathcal A\cup\{\eta E:E\in\mathcal B\}\);
the corresponding candidates in \eqref{eq:ordinaryfamily} have
the same ordinary answers.

A query containing both \(\xi\) and \(\eta\) is constantly positive
on \(\mathcal J\), so it is absent in the reduced tree.
The first nonordinary query therefore contains exactly one of them.
There are two cases.

\begin{enumerate}
\item If it contains \(\xi\), enter its positive child without testing.
From that point represent the possible inputs by
\[
 D\in\mathcal F\longmapsto \xi D,\qquad
 \eta E\in\mathcal C\longmapsto\xi\eta E.
\]
Every representative satisfies the previous ordinary answers and
contains \(\xi\), so it satisfies the skipped positive answer.
Later queries containing \(\xi\) are skipped into their positive
children. A query avoiding \(\xi\) is asked unchanged on
\(U\cup\{\eta\}\). Its OR answer equals the representative's answer,
because the only added defective label is \(\xi\).

\item If the first query contains \(\eta\), instead use
\[
 D\in\mathcal F\longmapsto\eta D,\qquad
 \eta E\in\mathcal C\longmapsto\xi\eta E.
\]
Again skip that query into its positive child, and skip any later
query containing \(\eta\). For a later query avoiding \(\eta\),
keep its ordinary labels and replace \(\xi\), if present, by the
label \(\eta\) of the \(\mathcal C\)-problem.
To check this replacement, write the old query as \(S\) or
\(S\cup\{\xi\}\), where \(S\subseteq U\). For a new input
\(D\in\mathcal F\), its representative is \(\eta D\):
\(\eta\) is absent from the new input and \(\xi\) is absent from
the representative. For a new input \(\eta E\), its representative
is \(\xi\eta E\): both of these compared labels are present.
In either case the ordinary part \(S\) meets precisely the same
ordinary defectives. Replacing \(\xi\) by \(\eta\) therefore preserves
the OR answer, including when \(S\) is empty.
\end{enumerate}

These representatives are part of a mathematical simulation, not
information revealed about the actual input. The choices of which
case applies depend only on the query reached in the original tree.
All asked questions are subsets of \(U\cup\{\eta\}\), and are
ordinary pooled-OR tests. The representation in either case is
injective, so the identifying output of the old tree uniquely
determines the input of the new one. The original family has the
known count \(k+1\), the simulated family the known count \(k\);
each specified representation adds exactly one defective.

\paragraph{Check every path length.}
Any simulated path that reaches a distinguished query skips at
least that first query, so has length at most \(s-1\).
If no such query occurs, it ends at an ordinary-only leaf.
That leaf cannot identify \(\xi D\) or \(\eta D\), which are
indistinguishable by ordinary queries. It identifies \(\xi\eta E\),
and its depth is at most \(s-1\) by the parent argument above.
Its corresponding \(\mathcal C\)-input is uniquely \(\eta E\).
Thus every simulated path has length at most \(s-1\).
We have \(\Comp(\mathcal C)\le s-1\), completing the lower bound
and the proof of Lemma~\ref{lem:two}.\quad\hfill\(\square\)

When \(k=1\), the family \(\mathcal G\) consists of the empty set.
The same proof works: its extension argument requires no pair of
distinct members in that case, and a nonempty \(\mathcal B\)
already has just one member.

\subsection{The recurrence and its direction}

\begin{proposition}\label{prop:recurrence}
For integers \(r\ge2\) and \(N\ge r+2\),
\begin{equation}\label{eq:recurrence}
 \boxed{M(r,N)\ge
 \min\{1+M(r,N-1),\ 2+M(r-1,N-1)\}.}
\end{equation}
\end{proposition}
\begin{proof}
Take a reduced optimal tree for the complete family of \(r\)-subsets
of an \(N\)-element set. Its root cannot be a constant query.

If its pool is a singleton, the negative child is a complete
\((r,N-1)\) problem. The root plus that child gives height at least
\(1+M(r,N-1)\). The positive child may be harder, but ignoring it
cannot invalidate this lower bound on the maximum of child depths.

If its pool has at least two labels, choose distinct \(\xi,\eta\)
inside it. Every \(r\)-subset meeting this pair gives a positive
root answer, irrespective of the other labels in the pool.
Thus the positive child contains the family \(\mathcal J\) of
Lemma~\ref{lem:two}, with \(U\) the other \(N-2\) labels and
\(k=r-1\). The conditions \(r\ge2,N\ge r+2\) ensure
\(1\le k<|U|\).
By restriction and the identity, the child's complexity is at least
\(1+M(r-1,N-1)\). Add the root to get
\(2+M(r-1,N-1)\).

Every nonconstant root belongs to one of these cases. Each tree is
at least as deep as the bound for its case, and hence at least as
deep as the smaller of the two bounds. This explains the minimum
in \eqref{eq:recurrence}.
\end{proof}

The argument did not restrict the root to singleton or pair pools:
the second case includes every pool with two or more labels.
Nor did it assume a special structure for later queries.
Both possible outcomes remain part of a legal tree, even though
one child can suffice to certify a lower bound.

\section{Turning a supposed saving into a contradiction}\label{sec:finalproof}

We now combine the two parts of the proof.
The recurrence will carry a supposed improvement over individual
testing to a smaller complete-family problem near density \(1/4\).
The information inequality will say that this smaller problem cannot
have the amount of saving that the recurrence requires.

\subsection{Measure the saving}

For \(0<r<N\), define
\[
 \Delta(r,N)=N-1-M(r,N).
\]
This is a nonnegative integer by individual testing.
It is the number of tests the optimal worst-case bound saves compared
with \(N-1\).
A counterexample to the desired equality has \(\Delta(r,N)\ge1\).

Rearrange the recurrence carefully. The two expressions on its
right are
\[
 1+M(r,N-1)=N-1-\Delta(r,N-1)
\]
and
\[
 2+M(r-1,N-1)=N-\Delta(r-1,N-1).
\]
Subtracting their minimum from \(N-1\) gives
\begin{equation}\label{eq:savingrecurrence}
 \Delta(r,N)\le
 \max\{\Delta(r,N-1),\ \Delta(r-1,N-1)-1\}
\end{equation}
when \(r\ge2,N\ge r+2\).

Suppose the current saving is at least \(s_0\).
The maximum on the right must be at least \(s_0\).
Therefore at least one of the following is true:
\begin{center}
\begin{tabular}{lll}
\toprule
Move to&Change of parameters&Required saving there\\
\midrule
\((r,N-1)\)&Remove one from the size only&At least \(s_0\)\\
\((r-1,N-1)\)&Remove one from rank and size&At least \(s_0+1\)\\
\bottomrule
\end{tabular}
\end{center}
Here rank means the defective count. A move lowering it increases
the forced saving by one.

This sequence of parameter pairs is a mathematical consequence of
lower bounds. It is not an actual test execution removing known
items for free. We are not claiming to know which items to delete
from an unknown input. The recurrence asserts that one of two
separate, completely specified testing problems must have the
indicated complexity.

\subsection{Start from an alleged improvement}

Assume, for a contradiction, that
\[
 n\ge2^{31},\qquad d/n\ge\cstar,\qquad \Delta(d,n)\ge1.
\]
Use \eqref{eq:savingrecurrence} to choose an available move of the
indicated kind at each stage. After \(j\) moves lowering rank and
size, and \(k_{\mathrm{size}}\) moves lowering size only, the
current pair satisfies
\begin{equation}\label{eq:tracking}
 r=d-j,\qquad N=n-j-k_{\mathrm{size}},\qquad
 \Delta(r,N)\ge1+j.
\end{equation}
This is proved by induction on moves: size-only moves preserve
the saving lower bound; rank-lowering moves increase it by one.

Stop at the first pair for which \(p=r/N\le1/4\).
The starting density is above \(1/4\), since
\(\cstar>1/3>1/4\); the first inequality follows from
\(3\cdot369056>1000000\).

\subsection{Why the sequence must stop at the desired density}

The recurrence is unavailable only if \(r=1\) or \(N=r+1\).
While the density is greater than \(1/4\), the case \(r=1\) means
\(N=2\) or \(3\). The other unavailable case is the one-good-item
problem \(N=r+1\).
All these states have \(\Delta=0\) by Lemma~\ref{lem:endpoints},
contradicting the positive lower bound in \eqref{eq:tracking}.
Thus an unavailable state cannot end our sequence above \(1/4\).

Each move decreases \(N\) by one, so the sequence cannot continue
indefinitely. Every permitted successor is admissible:
the recurrence guards \(r\ge2,N\ge r+2\) ensure positive new rank
and new rank smaller than new size.
It follows that the specified stopping pair exists.

A size-only move raises the ratio \(r/N\). Therefore the first
crossing below \(1/4\) must be a rank-lowering move. Its predecessor
was \((r+1,N+1)\), with density strictly above \(1/4\).
From
\[
 \frac{r+1}{N+1}>\frac14
 \quad\Longrightarrow\quad 4r>N-3
\]
we obtain the useful rounding control
\begin{equation}\label{eq:stoppingdensity}
 \frac14-\frac{3}{4N}<p=\frac rN\le\frac14.
\end{equation}
There is no need to assume that \(N\) is divisible by four.
Equality \(p=1/4\) is allowed at the stopping state.

\subsection{The remaining instance is still large}

Since \(M(r,N)\ge0\), we have \(\Delta(r,N)\le N-1\).
Combining this with \eqref{eq:tracking} gives \(j\le N-2\).
At the stopping state \(r\le N/4\), so
\[
 d=r+j\le\frac N4+N-2=\frac54N-2.
\]
Consequently
\begin{equation}\label{eq:remaining-size}
 N\ge\frac45(d+2)>\frac n4\ge2^{29}.
\end{equation}
For the strict middle inequality, use \(d\ge\cstar n\) and
\(\cstar>5/16\), since
\(16\cdot369056=5904896>5000000\).
Thus \((4/5)\cstar>1/4\).

Because \(N>2^{29}>750\), the rounding term \(3/(4N)\) is less
than \(1/1000\). Equation~\eqref{eq:stoppingdensity} therefore
places the stopped density in
\[
 \frac{249}{1000}<p\le\frac14.
\]
This is inside the interval of the whole-tree inequality.
There are no small residual instances left to check.

\subsection{How much saving is forced at the stopped state?}

Let \(c=d/n\) be the original density. This is an arbitrary number
at least \(\cstar\), not necessarily equal to it.
Since \(k_{\mathrm{size}}\ge0\),
\[
 N\le n-j=n-d+r=n-d+pN.
\]
Rearranging, with \(c<1\), gives
\[
 \frac nN\ge\frac{1-p}{1-c}.
\]
Since \(j=d-r\),
\begin{equation}\label{eq:rankmoves}
 \frac jN=\frac dN-p
 \ge\frac{c(1-p)}{1-c}-p
 =\frac{c-p}{1-c}.
\end{equation}
For fixed \(p<1\), the derivative of the last expression with
respect to \(c\) is \((1-p)/(1-c)^2>0\).
For fixed \(c<1\), its derivative with respect to \(p\) is
\(-1/(1-c)<0\).
Using \(c\ge\cstar\) and \(p\le1/4\), we conclude
\[
 \frac jN\ge
 \frac{\cstar-1/4}{1-\cstar}
 =\frac{119056}{630944}
 =\frac{7441}{39434}.
\]
Define
\begin{equation}\label{eq:gstar}
 g_\star=\frac{7441}{39434}.
\end{equation}
The tracked saving is therefore at least
\begin{equation}\label{eq:lower-saving}
 \Delta(r,N)\ge1+g_\star N.
\end{equation}
This is why all starting densities above \(\cstar\) are covered:
increasing the original density only strengthens the lower bound
on \(j/N\).

\subsection{What the information bound allows}

Apply Theorem~\ref{thm:whole} to an optimal tree for the stopped
\((r,N)\) problem. Here use
\[
 H=\log\binom Nr,\qquad K=Nh(p)-H.
\]
The theorem and expected depth at most worst-case depth give
\[
 M(r,N)\ge H+\frac{N-H-67229K}{6615}.
\]
Subtract from \(N-1\), and show each coefficient:
\begin{align*}
 \Delta(r,N)
 &\le N-1-H-\frac{N-H}{6615}+\frac{67229}{6615}K\\
 &=\frac{6614}{6615}(N-H)+\frac{67229}{6615}K-1\\
 &=\frac{6614}{6615}N(1-h(p))
    +\frac{6614+67229}{6615}K-1.
\end{align*}
Put
\[
 \alpha=\frac{6614}{6615},\qquad
 \beta=\frac{73843}{6615}<12.
\]
The last comparison follows from \(73843<12\cdot6615=79380\).
Using \(K\le\log(N+1)\), we obtain
\begin{equation}\label{eq:upper-saving}
 \Delta(r,N)\le
 \alpha N(1-h(p))+\beta\log(N+1)-1.
\end{equation}

We next replace the nearby density \(p\) by exactly \(1/4\), paying
only a constant for the rounding. Define
\[
 \er=1-h(1/4).
\]
For \(249/1000\le u\le1/4\),
\[
 h'(u)=\log\frac{1-u}{u}<2,
\]
since \((1-u)/u\le751/249<4\).
Integration from \(p\) to \(1/4\), together with
\eqref{eq:stoppingdensity}, yields
\[
 h(1/4)-h(p)\le2(1/4-p)<\frac3{2N}.
\]
It also holds when \(p=1/4\), because then the left side is zero.
Thus
\begin{equation}\label{eq:roundedentropy}
 N(1-h(p))<N\er+\frac32.
\end{equation}

\subsection{The numerical margin and the contradiction}

The exact arithmetic needed now is isolated in the next lemma.
Its proof, including finite logarithm series and their remainder
bounds, is given in Appendix~\ref{app:arithmetic}.

\begin{lemma}[Two numerical certificates]\label{lem:numerics}
With the constants above,
\begin{equation}\label{eq:positive-margin}
 g_\star-\alpha\er>\frac1{1000000}.
\end{equation}
For every real \(N\ge2^{29}\),
\begin{equation}\label{eq:logcutoff}
 12\log(N+1)<\frac N{1000000}.
\end{equation}
\end{lemma}

Combine the lower saving \eqref{eq:lower-saving}, the upper saving
\eqref{eq:upper-saving}, and the rounded entropy estimate:
\[
 1+g_\star N
 \le\Delta(r,N)
 <\alpha N\er+\frac{3\alpha}{2}
   +\beta\log(N+1)-1.
\]
Move the linear terms to the left and the two constant ones to
the right:
\[
 (g_\star-\alpha\er)N
 <\beta\log(N+1)+\frac{3\alpha}{2}-2.
\]
Since \(0<\alpha<1\), the final constant is negative. Hence a
counterexample would require
\begin{equation}\label{eq:contradiction}
 (g_\star-\alpha\er)N<\beta\log(N+1).
\end{equation}
But Lemma~\ref{lem:numerics}, \(\beta<12\), and
\(N>2^{29}\) give the incompatible chain
\[
 (g_\star-\alpha\er)N
 >\frac N{1000000}
 >12\log(N+1)
 >\beta\log(N+1).
\]
This contradiction rules out \(\Delta(d,n)\ge1\).
As \(\Delta(d,n)\) is a nonnegative integer, it must equal zero:
\[
 M(d,n)=n-1.
\]
Together with Proposition~\ref{prop:individual}, this proves
Theorem~\ref{thm:main}.\quad\hfill\(\square\)

\subsection{What each constant does}

The proof separates several roles that are easy to confuse.
\begin{itemize}
\item The intermediate density \(1/4\) makes large pools and fresh
pairs biased in useful directions under their respective comparison
laws. The small interval around it absorbs the integer crossing error.
\item The coefficient \(6615\) comes from the tree accounting. It
produces \(\alpha=1-1/6615<1\) in the upper bound on savings.
\item The starting threshold \(\cstar\) determines
\(g_\star=(\cstar-1/4)/(1-\cstar)\), the saving forced per remaining
item by the recurrence.
\item The margin \(g_\star-\alpha\er>10^{-6}\) means that the forced
linear saving exceeds what the information estimate permits, apart
from its error term.
\item The exact-count comparison contributes \(K\le\log(N+1)\),
giving a logarithmic error. The cutoff \(n\ge2^{31}\) guarantees a
residual \(N>2^{29}\), where the linear margin dominates that error.
\end{itemize}
The numerical certificate is used to establish these inequalities,
not to guess a pattern from small instances.

\section{What has, and has not, been proved}\label{sec:scope}

The conclusion is an exact worst-case equality, for every admissible
integer pair in the stated range. The proof allowed arbitrary legal
pools, arbitrary overlap, and full adaptivity. The average used in
the probability argument was always over valid exact-count inputs.
The independent and artificial laws were comparison measures, with
every change of law accounted for by an identity or inequality.

The expectation estimate alone was not the final result. It supplied
an upper bound on possible savings at a smaller instance. The
structural recurrence forced an incompatible lower bound on those
savings from any proposed original counterexample. This is the step
that upgrades the information estimate to the exact answer \(n-1\).

There is no assertion that the same threshold theorem holds for all
\(n<2^{31}\). The elementary endpoint results proved earlier do hold
at their stated smaller sizes, but they are not a substitute for a
general small-population theorem.
The constants in the argument are sufficient constants; their
optimality is not claimed.

The conjectural boundary is density \(1/3\), often called the
Hu--Hwang--Wang conjecture \cite{HHW}. Since
\(3\cdot369056>1000000\), the theorem explained here stops strictly
above that boundary. No resolution at \(1/3\) follows from this proof.

\appendix
\section{Exact certificates for the numerical step}\label{app:arithmetic}

This appendix proves Lemma~\ref{lem:numerics}. Its purpose is narrow:
to certify a positive linear margin and the population size at which
that margin dominates the logarithmic correction. Decimal
approximations are not used as evidence.

\subsection{A logarithm series with a proved remainder}

For \(0<t<1\), the derivative of
\(\ln((1+t)/(1-t))\) is \(2/(1-t^2)\).
The finite geometric identity gives
\[
 \frac1{1-u^2}=\sum_{j=0}^{s}u^{2j}
                +\frac{u^{2s+2}}{1-u^2}.
\]
Integrate from zero to \(t\), multiply by two, and use that the
logarithm is zero at zero:
\begin{equation}\label{eq:logseries}
 \ln\frac{1+t}{1-t}
 =2\sum_{j=0}^{s}\frac{t^{2j+1}}{2j+1}+R_s(t).
\end{equation}
The remainder is the integral of \(2u^{2s+2}/(1-u^2)\).
It is positive, and \(1-u^2>1-t^2\) for \(0\le u<t\).
Therefore
\begin{equation}\label{eq:logremainder}
 0<R_s(t)<\frac{2t^{2s+3}}{(2s+3)(1-t^2)}.
\end{equation}
At \(t=1/3\) the logarithm in \eqref{eq:logseries} is \(\ln2\);
at \(t=1/5\) it is \(\ln(3/2)\).

\subsection{A finite rational table}

Set \(T_0=10^{10}\).
Every row in the following table has positive integers \(u,f\)
satisfying
\[
 uf\le2T_0<u(f+1).
\]
Equivalently, \(f/T_0\le2/u<(f+1)/T_0\).
Thus the table gives lower and upper rational bounds for each
series term; it is a certificate consisting only of integer
multiplications and comparisons.

\begin{center}
\begin{tabular}{rrr@{\qquad}rrr}
\toprule
\multicolumn{3}{c}{For \(\ln2\), \(t=1/3\)}&
\multicolumn{3}{c}{For \(\ln(3/2)\), \(t=1/5\)}\\
\(j\)&\(u=(2j+1)3^{2j+1}\)&\(f\)&
\(j\)&\(u=(2j+1)5^{2j+1}\)&\(f\)\\
\midrule
0&3&6666666666&0&5&4000000000\\
1&81&246913580&1&375&53333333\\
2&1215&16460905&2&15625&1280000\\
3&15309&1306421&3&546875&36571\\
4&177147&112900&4&17578125&1137\\
5&1948617&10263&5&537109375&37\\
6&20726199&964&&&\\
7&215233605&92&&&\\
8&2195382771&9&&&\\
\midrule
\multicolumn{2}{r}{Sum of \(f\):}&6931471800&
\multicolumn{2}{r}{Sum of \(f\):}&4054651078\\
\bottomrule
\end{tabular}
\end{center}

For example, in the row \(u=81\) the product is
\(81\cdot246913580=19999999980\), while adding another \(81\)
exceeds \(20000000000=2T_0\).
The displayed inequalities for the other rows have exactly the
same form.

The remainder bounds \eqref{eq:logremainder} give
\[
 R_8(1/3)<\frac1{76\cdot3^{17}}=\frac1{9814652388}
 <\frac2{T_0},
\]
and
\[
 R_5(1/5)<\frac1{156\cdot5^{11}}=\frac1{7617187500}
 <\frac2{T_0}.
\]
In the first sum the error from replacing the nine terms by their
table bounds is less than \(9/T_0\). In the second it is less than
\(6/T_0\). Hence
\[
 \frac{6931471800}{T_0}<\ln2<
 \frac{6931471800+9+2}{T_0}
 <\frac{6931471900}{T_0},
\]
and
\[
 \frac{4054651000}{T_0}
 <\frac{4054651078}{T_0}<\ln(3/2)
 <\frac{4054651078+6+2}{T_0}
 <\frac{4054651100}{T_0}.
\]
In particular,
\begin{equation}\label{eq:lncert}
 \frac{69314718}{10^8}<\ln2<\frac{69314719}{10^8},
 \qquad
 \frac{40546510}{10^8}<\ln(3/2)<\frac{40546511}{10^8}.
\end{equation}
These also certify the simpler bound \(\ln2>1/2\).

\subsection{Bound the entropy gap at one quarter}

Let \(\theta=\log_2 3\). Since
\(\theta=1+\ln(3/2)/\ln2\), the upper bound on the numerator and
the lower bound on the denominator in \eqref{eq:lncert} yield
\[
 \theta<1+\frac{40546511}{69314718}
 <\frac{1584963}{1000000}.
\]
The last comparison has the exact integer certificate
\[
 584963\cdot69314718-1000000\cdot40546511
 =34385434>0.
\]
From \(h(1/4)=2-3\theta/4\), we obtain
\begin{equation}\label{eq:ebar}
 \er=1-h(1/4)=\frac34\theta-1
 <\bar e:=\frac{754889}{4000000}.
\end{equation}

\subsection{The positive margin}

Recall \(g_\star=7441/39434\) and \(\alpha=6614/6615\).
Exact subtraction gives
\[
 g_\star-\bar e
 =-\frac{4292826}{157736000000}>
 -\frac{55}{2000000}.
\]
After multiplication by the positive denominator, the final
inequality is \(4292826<4337740\).
Also
\[
 \frac{\bar e}{6615}
 =\frac{754889}{26460000000}>
 \frac{57}{2000000},
\]
whose certificate is \(754889>754110\).
Now combine these as parts of one expression:
\begin{align*}
 g_\star-\alpha\er
 &>g_\star-\alpha\bar e\\
 &=(g_\star-\bar e)+\frac{\bar e}{6615}\\
 &>-\frac{55}{2000000}+\frac{57}{2000000}
 =\frac1{1000000}.
\end{align*}
The first inequality uses \(\alpha>0\) and \(\er<\bar e\).
This proves \eqref{eq:positive-margin}.

\subsection{The cutoff}

For real \(N\ge2^{29}\), define
\[
 F(N)=\frac N{1000000}-12\log(N+1).
\]
At \(N=2^{29}=536870912\), the first term exceeds \(536\).
Since \(2^{29}+1<2^{30}\), the second term is less than
\(12\cdot30=360\). Thus \(F(2^{29})>0\).
Furthermore, using \(\ln2>1/2\),
\[
 F'(N)=\frac1{1000000}-\frac{12}{(N+1)\ln2}
 >\frac1{1000000}-\frac{24}{N+1}>0.
\]
The last inequality holds because
\(N+1\ge2^{29}+1>24000000\).
So \(F(N)>0\) for the entire range, proving
\eqref{eq:logcutoff}.
This completes the proof of Lemma~\ref{lem:numerics}.

\section{Notation guide}\label{app:notation}

The following guide collects the main notation; definitions and
justifications are in the body of the exposition.
\begin{longtable}{p{0.23\linewidth}p{0.70\linewidth}}
\toprule
Symbol&Meaning\\
\midrule
\endfirsthead
\toprule
Symbol&Meaning\\
\midrule
\endhead
\bottomrule
\endfoot
\(n,d\)&Original population size and exact defective count.\\
\(D\)&A particular defective subset.\\
\(X,\ X_i\)&Random defective subset and its indicator at label \(i\).\\
\(S_v,\widehat S_v\)&Original and normalized pools at node \(v\).\\
\(Z_v\)&Union of preceding pools assigned negative answers on the
node's fixed prefix.\\
\(\tau,\ y,\ \ell\)&Prefix transcript, terminal transcript, and terminal leaf.\\
\(Y,T\)&Random terminal transcript and number of tests in an execution.\\
\(P_X,P_Y,P_{X,Y}\)&Input, transcript, and joint laws for uniform
exact-count inputs with truthful answers.\\
\(Q_X,Q_Y,Q_{X,Y}\)&Corresponding laws for independent inputs with truthful answers.\\
\(R_X,R_Y,R_{X,Y}\)&Corresponding laws for independent inputs with artificial
fair answers at large normalized pools. Here \(R_X=Q_X\).\\
\(P(v)\)&Probability under \(P\) of reaching node \(v\).\\
\(\pi_v,\rho_v^Q,\rho_v^R\)&Conditional negative probabilities at \(v\)
under \(P,Q,R\), respectively.\\
\(\Ent,\ h,\ \KL,\ \kl\)&Finite-law entropy, binary entropy, relative entropy,
and binary relative entropy. All use base-two logarithms.\\
\(p,q\)&The current exact defective fraction and its complement \(q=1-p\).\\
\(H\)&\(\log\binom nd\), entropy of the uniform valid-input law.\\
\(K\)&\(nh(p)-H=\KL(P_X\Vert Q_X)\), the exact-count comparison cost.\\
\(L,\Gamma\)&Actual expected test count \(\E_P T\) and its excess \(L-H\).\\
\(\Glg,\Gfr\)&Parts of \(\Gamma\) contributed by large queries and fresh pairs.\\
\(a(D),b(D),t(D)\)&Large-query count, fresh-pair count, and test count
on the truthful path of \(D\).\\
\(A,B\)&The \(P\)-expected large-query and fresh-pair counts.\\
\(V(D),w_1(D),w_0(D)\)&Union of small pools on the truthful path, and
counts of defective and good labels outside that union.\\
\(W\)&\(\E_P(w_1+w_0)\), expected number of untouched labels.\\
\(J\)&\(\E_P[w_1\log(1/p)+w_0\log(1/q)]\), a lower bound on averaged
conditional input divergence.\\
\(\mathcal D_R\)&\(\KL(P_Y\Vert R_Y)\), reference transcript divergence.\\
\(\Comp(\mathcal F)\)&Minimum worst-case depth for a specified nonempty
candidate family \(\mathcal F\).\\
\(\xi,\eta\)&The two distinguished labels in the structural identity.\\
\(r,N\)&Rank and population size at a smaller recurrence instance.\\
\(\Delta(r,N)\)&\(N-1-M(r,N)\), saving compared with individual testing.\\
\(j,k_{\mathrm{size}}\)&Counts of rank-and-size moves and size-only moves
in the recurrence sequence.\\
\(\cstar,g_\star\)&\(369056/1000000\) and
\((\cstar-1/4)/(1-\cstar)=7441/39434\).\\
\(\alpha,\beta\)&\(6614/6615\) and \(73843/6615\), coefficients in the
upper bound on savings.\\
\(\er,\bar e\)&The gap \(1-h(1/4)\) and its rational upper bound
\(754889/4000000\).\\
\end{longtable}

\begin{thebibliography}{9}
\bibitem{Source}
\emph{Coupled Reference-Process Bounds for Dense Adaptive Group Testing},
unpublished mathematical manuscript, October 25, 2026.
Mathematical source for this exposition; all results needed here
have been proved within the present document.

\bibitem{HHW}
M.~C. Hu, F.~K. Hwang, and J.~K. Wang,
A boundary problem for group testing,
\emph{SIAM Journal on Algebraic and Discrete Methods}
\textbf{2}(2) (1981), 81--87.
\href{https://doi.org/10.1137/0602011}{doi:10.1137/0602011}.
\end{thebibliography}
\end{document}
